\section{Datos espectrales de una realización supersimétrica: \((\mathcal H_{\rm phys},H,\mathcal Q)\), supercarga y compactificación en dimensión \(11\)}
\label{sec:md-sucesora-terna-supercarga-m}

La interfaz dimensional hacia teoría M exige distinguir la estructura exacta
de partida, el levantamiento operatorio y los campos físicos del codominio. Sea
\[
 \mathfrak S_M=(\mathcal H_{\mathrm{phys}},H,Q)
\]
una terna espectral en la que \(\mathcal H_{\mathrm{phys}}\) es el espacio de
Hilbert después de imponer las restricciones de calibre, \(H\) es un operador
autoadjunto no negativo y \(Q\) es una supercarga impar, densamente definida y
compatible con la graduación
\(\mathcal H_{\mathrm{phys}}=\mathcal H_+\oplus\mathcal H_-\). La relación
\begin{equation}
 Q:\mathcal H_+\rightleftarrows\mathcal H_-,
 \qquad
 Q^2=H
 \label{eq:md-sucesora-supercarga-cuadrado}
\end{equation}
se interpreta en su dominio común. No reemplaza el álgebra supersimétrica
relativista; es su reducción hermítica de dimensión \(1\) y fija el dato mínimo
que debe conservar una realización espectral.

\subsection{Levantamiento operatorio del estado APP--TRIT--TPK}

Sea \(\mathcal H_{\mathrm{HMT}}\) la completación de las secciones
supervivientes con fase, orientación, cociente, acarreo, frontera, ruta y
memoria, y sea
\[
 W:\mathcal H_{\mathrm{HMT}}\longrightarrow\mathcal H_{\mathrm{phys}}
\]
una isometría de realización. Un levantamiento operatorio admisible conserva
la composición de APP, la orientación de TRIT y la dinámica completa del TPK,
y entrelaza
\begin{equation}
 W\widehat H=HW,
 \qquad
 W\widehat Q=QW,
 \qquad
 \widehat Q^{2}=\widehat H.
 \label{eq:md-sucesora-lift-app-supercarga}
\end{equation}
La primera igualdad transporta el generador, la segunda preserva la graduación
y la tercera impide introducir \(H\) como un dato ajeno al antecedente
genealógico. Estas identidades son obligaciones del puente: la existencia de
completaciones fermiónica y bosónica, por sí sola, no demuestra
\eqref{eq:md-sucesora-supercarga-cuadrado}.

\Needspace{42\baselineskip}
\subsection{Codominio físico undecadimensional}

La realización en dimensión \(11\) se expresa mediante
\[
 \mathfrak R_M:\mathcal X_{\mathrm{HMT}}^{\mathrm{enr}}
 \longrightarrow
 \mathcal X_{11}^{\mathrm{phys}},
\]
cuyo codominio porta la métrica \(g_{MN}\), la \(3\)-forma \(C_{MNP}\), el
gravitino \(\psi_M\), la acción y sus transformaciones de supersimetría. La
compactificación sobre una dirección debe transportar simultáneamente estos
campos, el producto interno y la graduación de \(\mathfrak S_M\). Las cuerdas
aparecen en la pantalla de dimensión (10); las branas se acoplan a la
\(3\)-forma y a sus duales; y la dualidad T intercambia momento y
enrollamiento mediante
\[
 (m,w,R)\longmapsto(w,m,R^{-1})
\]
sin alterar el espectro del sector compacto.

La acción de Polyakov
\[
 S_{\mathrm P}[X,h]
 =-\frac{T}{2}\int d^2\sigma\,
 \sqrt{-h}\,h^{ab}g_{MN}(X)
 \partial_aX^M\partial_bX^N
\]
fija la dinámica de la superficie de mundo una vez determinados \(T\),
\(h_{ab}\), los mapas \(X^M\) y las condiciones de frontera. La cadena de
pantallas \(12\to11\to10\) y \(26\to24\), la involución de radio y la
reducción reticular son resultados algebraicos anteriores; la acción, el
espectro y las amplitudes constituyen la realización física posterior.

\subsection{Compatibilidad de la terna espectral \((\mathcal H_{\rm phys},H,\mathcal Q)\) con el codominio físico en dimensión \(11\)}

Una realización es admisible si preserva a la vez:
\begin{enumerate}
 \item la prolongación compatible del estado a toda profundidad;
 \item los dominios y codominios de \(H\), \(Q\) y los observables;
 \item las identidades de entrelazamiento
       \eqref{eq:md-sucesora-lift-app-supercarga};
 \item la acción unitaria de la dualidad T sobre el espectro;
 \item el descenso de los campos de dimensión \(11\) a la compactificación
       declarada;
 \item la compatibilidad con el producto interno físico y las restricciones
       de calibre.
\end{enumerate}
El fallo de cualquiera de estas condiciones distingue una analogía de rangos
de una realización física. La formulación no identifica automáticamente el
corredor excepcional con teoría M: ambas publicaciones proceden del mismo
estado enriquecido, pero utilizan funtores y codominios distintos.
